\documentclass[11pt]{article}
\usepackage[margin=1in]{geometry}
\usepackage[T1]{fontenc}
\usepackage[utf8]{inputenc}
\usepackage[french]{babel}
\usepackage{amsmath,amssymb,amsthm}
\usepackage[colorlinks=true,linkcolor=blue,urlcolor=blue]{hyperref}
\newtheorem{problem}{Problème}
\newenvironment{refs}{\par\small\noindent\emph{Références :}\begin{itemize}\setlength\itemsep{0pt}}{\end{itemize}\normalsize}
\title{Hors de la Caverne \\ \Large Théorie des catégories}
\author{}
\date{}
\begin{document}
\maketitle
\noindent\emph{Des problèmes, pas de réponses.} Chaque problème ci-dessous est posé
sans solution. Les références indiquent où trouver les connaissances nécessaires.
\medskip


\section*{Licence}

\begin{problem}[{Est-ce une catégorie ? --- Licence}]
\textbf{CAT-01.} \textbf{Problème.} Une catégorie est constituée d'objets, de morphismes munis d'une
source et d'un but, d'une composition associative, et d'un morphisme identité sur chaque objet.
\begin{enumerate}
\item[(a)] Vérifiez les axiomes pour : les ensembles avec les applications ($ \mathbf{Set} $) ; les
groupes avec les homomorphismes ($ \mathbf{Grp} $) ; un ensemble partiellement ordonné
$ (P, \le) $ avec exactement un morphisme $ p \to q $ lorsque $ p \le q $ et aucun sinon ;
un monoïde unique vu comme une catégorie à un seul objet.
\item[(b)] Démontrez que le morphisme identité d'un objet est unique.
\item[(c)] On appelle
$ f\colon A \to B $ un \textit{monomorphisme} si $ f \circ g = f \circ h $ entraîne
$ g = h $, et un \textit{épimorphisme} si $ g \circ f = h \circ f $ entraîne $ g = h $.
Démontrez que dans $ \mathbf{Set} $ les monomorphismes sont exactement les injections et les
épimorphismes exactement les surjections.
\item[(d)] Montrez que dans la catégorie des anneaux l'inclusion
$ \mathbb{Z} \hookrightarrow \mathbb{Q} $ est un épimorphisme qui n'est pas surjectif.
\end{enumerate}

\end{problem}

\noindent Les catégories ont été introduites par Samuel Eilenberg et Saunders Mac Lane
dans leur article de 1945 \og{} General theory of natural equivalences \fg{} --- le nom emprunté à Kant,
comme \og{} foncteur \fg{} l'était à Carnap. Ils insistaient sur le fait que ces définitions n'étaient
qu'une comptabilité au service du véritable but : la naturalité (CAT-05). La partie (d) est
l'avertissement classique que les notions définies par les flèches peuvent diverger de leurs
modèles ensemblistes : \og{} surjectif \fg{} n'est pas un concept que toute catégorie sait voir.

\begin{refs}
\item Contexte : Théorie des catégories --- Wikipédia (\url{https://fr.wikipedia.org/wiki/Th%C3%A9orie_des_cat%C3%A9gories})
\item Contexte : Monomorphisme --- Wikipédia (\url{https://fr.wikipedia.org/wiki/Monomorphisme}) et Épimorphisme --- Wikipédia (\url{https://fr.wikipedia.org/wiki/%C3%89pimorphisme})
\item Lecture : Tom Leinster, \textit{Basic Category Theory}, ch. 1 (librement accessible sur arXiv:1612.09375 (\url{https://arxiv.org/abs/1612.09375}))
\item Article : S. Eilenberg et S. Mac Lane, \og{} General theory of natural equivalences \fg{} (1945), Trans. Amer. Math. Soc. 58
\end{refs}

\bigskip

\begin{problem}[{Ce que préservent les foncteurs --- Licence}]
\textbf{CAT-02.} \textbf{Problème.} Un foncteur $ F\colon \mathcal{C} \to \mathcal{D} $
associe des objets à des objets et des morphismes à des morphismes, en préservant sources, buts,
composition et identités.
\begin{enumerate}
\item[(a)] Vérifiez que chacune des constructions suivantes est un foncteur : le foncteur d'oubli
$ U\colon \mathbf{Grp} \to \mathbf{Set} $ ; le foncteur parties
$ \mathcal{P}\colon \mathbf{Set} \to \mathbf{Set} $ avec $ \mathcal{P}(f)(S) = f[S] $ ;
l'abélianisation $ G \mapsto G/[G,G] $ de $ \mathbf{Grp} $ vers $ \mathbf{Ab} $.
\item[(b)] Démontrez que tout foncteur préserve les isomorphismes, et concluez-en le principe
contraposé : si $ F(X) \not\cong F(Y) $ alors $ X \not\cong Y $.
\item[(c)] Démontrez qu'il n'existe
\textit{aucun} foncteur $ \mathbf{Grp} \to \mathbf{Grp} $ envoyant chaque groupe sur son centre.
\end{enumerate}

\end{problem}

\noindent La partie (b) est la licence qui autorise tout invariant en mathématiques --- le
groupe fondamental, l'homologie, les déterminants : calculez l'image, distinguez les objets. La
partie (c), exercice favori du livre de Riehl, est la leçon la plus tranchante qui soit sur le
fait qu'une \og{} construction sur les objets \fg{} n'est pas encore une construction : la fonctorialité
est une contrainte véritable, et le centre --- si naturel qu'il paraisse --- ne la satisfait pas.

\begin{refs}
\item Contexte : Foncteur --- Wikipédia (\url{https://fr.wikipedia.org/wiki/Foncteur})
\item Contexte : Foncteur d'oubli --- Wikipédia (en anglais) (\url{https://en.wikipedia.org/wiki/Forgetful_functor})
\item Lecture : Emily Riehl, \textit{Category Theory in Context}, ch. 1 (librement accessible sur emilyriehl.github.io/files/context.pdf (\url{https://emilyriehl.github.io/files/context.pdf}))
\end{refs}

\bigskip

\begin{problem}[{Produits et coproduits comme propriétés universelles --- Licence}]
\textbf{CAT-03.} \textbf{Problème.} Un \textit{produit} des objets $ A $ et $ B $ est un
objet $ P $ muni de morphismes $ \pi_1\colon P \to A $, $ \pi_2\colon P \to B $ tels que
pour tout objet $ X $ muni de morphismes $ f\colon X \to A $, $ g\colon X \to B $ il existe
un unique $ h\colon X \to P $ avec $ \pi_1 \circ h = f $ et $ \pi_2 \circ h = g $. Un
\textit{coproduit} est la notion duale, toutes les flèches étant renversées.
\begin{enumerate}
\item[(a)] Démontrez que deux
produits quelconques de $ A $ et $ B $ sont isomorphes par un unique isomorphisme compatible
avec les projections.
\item[(b)] Identifiez, avec démonstration, le produit et le coproduit dans : $ \mathbf{Set} $
(produit cartésien ; réunion disjointe) ; $ \mathbf{Grp} $ (produit direct ; produit libre --- en
particulier, vérifiez que $ \mathbb{Z} * \mathbb{Z} $, le groupe libre à deux générateurs,
satisfait la propriété du coproduit) ; $ \mathbf{Ab} $ (produit direct ; somme directe) ; un
ensemble ordonné vu comme une catégorie (borne inférieure ; borne supérieure).
\item[(c)] Exhibez une catégorie dans laquelle un certain couple d'objets n'a
pas de produit, et démontrez-le.
\end{enumerate}

\end{problem}

\noindent Définir un objet par sa propriété universelle --- par la manière dont tout le
reste s'envoie vers lui, plutôt que par ce qu'il contient --- est le geste signature de la théorie
des catégories. L'article de Mac Lane \og{} Duality for groups \fg{} (1950) fut le premier à caractériser
ainsi les produits directs et à observer qu'en renversant les flèches on obtient le produit libre :
une définition, deux théorèmes. L'exemple des ensembles ordonnés montre le même schéma régissant
la théorie de l'ordre : les produits y sont des bornes inférieures.

\begin{refs}
\item Contexte : Produit (catégorie) --- Wikipédia (\url{https://fr.wikipedia.org/wiki/Produit_%28cat%C3%A9gorie%29})
\item Contexte : Somme, ou coproduit (catégorie) --- Wikipédia (\url{https://fr.wikipedia.org/wiki/Somme_%28cat%C3%A9gorie%29})
\item Contexte : Propriété universelle --- Wikipédia (\url{https://fr.wikipedia.org/wiki/Propri%C3%A9t%C3%A9_universelle})
\item Article : S. Mac Lane, \og{} Duality for groups \fg{} (1950), Bull. Amer. Math. Soc. 56
\end{refs}

\bigskip

\begin{problem}[{Objet initial, objet final, et la puissance de la dualité --- Licence}]
\textbf{CAT-04.} \textbf{Problème.} Un objet est \textit{initial} s'il existe exactement un
morphisme de lui vers chaque objet, et \textit{final} s'il existe exactement un morphisme vers lui
depuis chaque objet.
\begin{enumerate}
\item[(a)] Démontrez que deux objets initiaux quelconques sont isomorphes par un unique
isomorphisme.
\item[(b)] Identifiez, avec démonstration, les objets initiaux et finaux (ou montrez qu'il n'en existe
pas) dans : $ \mathbf{Set} $ ; $ \mathbf{Grp} $ ; la catégorie des anneaux unitaires ; la
catégorie des corps ; un ensemble ordonné vu comme une catégorie.
\item[(c)] La \textit{catégorie opposée}
$ \mathcal{C}^{\mathrm{op}} $ a les mêmes objets et les morphismes renversés. Expliquez
rigoureusement comment (a) fournit, sans travail supplémentaire, le théorème d'unicité
correspondant pour les objets finaux --- et pourquoi tout théorème catégorique arrive accompagné
d'un dual gratuit.
\end{enumerate}

\end{problem}

\noindent Le principe de dualité est l'offre \og{} deux pour le prix d'un \fg{} de la théorie des
catégories, héritée de la dualité entre points et droites en géométrie projective. Il explique un
air de famille qui traverse toutes les mathématiques : noyaux et conoyaux, limites et colimites,
produits et coproduits sont des théorèmes uniques portant deux costumes. La catégorie des corps,
qui n'a ni objet initial, ni objet final, ni produits, est le contre-exemple permanent qui
maintient le sujet honnête.

\begin{refs}
\item Contexte : Objet initial et objet final --- Wikipédia (\url{https://fr.wikipedia.org/wiki/Objet_initial_et_objet_final})
\item Contexte : Catégorie opposée --- Wikipédia (\url{https://fr.wikipedia.org/wiki/Cat%C3%A9gorie_oppos%C3%A9e})
\item Lecture : Steve Awodey, \textit{Category Theory}, ch. 3
\end{refs}

\bigskip

\begin{problem}[{Transformations naturelles --- Licence}]
\textbf{CAT-05.} \textbf{Problème.} Une \textit{transformation naturelle} $ \eta $ entre
foncteurs $ F, G\colon \mathcal{C} \to \mathcal{D} $ associe à chaque objet $ A $ un
morphisme $ \eta_A\colon F(A) \to G(A) $ tel que $ G(f) \circ \eta_A = \eta_B \circ F(f) $
pour tout $ f\colon A \to B $. Travaillez dans la catégorie des espaces vectoriels sur un corps $ k $.
\begin{enumerate}
\item[(a)] Démontrez que les applications d'évaluation $ \eta_V\colon V \to V^{**} $, $ \eta_V(v)(\varphi) =
\varphi(v) $, forment une transformation naturelle du foncteur identité vers le foncteur
bidual.
\item[(b)] La dualité simple $ V \mapsto V^{*} $ est contravariante, si bien que l'énoncé
\og{} isomorphisme naturel $ V \cong V^{*} $ \fg{} demande de la précaution. Sur la catégorie dont les
objets sont les espaces de dimension finie et dont les morphismes sont seulement les
isomorphismes, $ V \mapsto V^{*} $, $ f \mapsto (f^{-1})^{*} $ est un foncteur covariant.
Démontrez que pour $ k = \mathbb{R} $ toute transformation naturelle de l'identité vers ce
foncteur est identiquement nulle --- donc qu'aucun isomorphisme naturel $ V \cong V^{*} $
n'existe, même là.
\item[(c)] Démontrez que les foncteurs
$ \mathcal{C} \to \mathcal{D} $ et les transformations naturelles entre eux forment une
catégorie.
\end{enumerate}

\end{problem}

\noindent Ce problème est la raison d'être de la théorie des catégories : l'article de
1945 d'Eilenberg et Mac Lane fut écrit pour donner un foyer précis à l'intuition du praticien
selon laquelle $ V \cong V^{**} $ est canonique tandis que $ V \cong V^{*} $ exige un choix
arbitraire de base. Leur formule --- les catégories ont été définies pour pouvoir définir les
foncteurs, et les foncteurs pour pouvoir définir les transformations naturelles --- est
historiquement littérale.

\begin{refs}
\item Contexte : Transformation naturelle --- Wikipédia (\url{https://fr.wikipedia.org/wiki/Transformation_naturelle})
\item Lecture : Emily Riehl, \textit{Category Theory in Context}, ch. 1
\item Article : S. Eilenberg et S. Mac Lane, \og{} General theory of natural equivalences \fg{} (1945), Trans. Amer. Math. Soc. 58
\end{refs}

\bigskip

\begin{problem}[{Isomorphisme contre équivalence de catégories --- Licence}]
\textbf{CAT-06.} \textbf{Problème.} Les catégories $ \mathcal{C} $ et $ \mathcal{D} $ sont
\textit{isomorphes} s'il existe des foncteurs $ F\colon \mathcal{C} \to \mathcal{D} $ et
$ G\colon \mathcal{D} \to \mathcal{C} $ avec $ GF = \mathrm{id} $ et
$ FG = \mathrm{id} $ ; elles sont \textit{équivalentes} si l'on a plutôt $ GF \cong \mathrm{id} $ et
$ FG \cong \mathrm{id} $ via des isomorphismes naturels.
\begin{enumerate}
\item[(a)] Démontrez que $ F $ est une moitié d'une
équivalence si et seulement si $ F $ est plein, fidèle et essentiellement surjectif (tout
objet de $ \mathcal{D} $ est isomorphe à un $ F(A) $) --- et repérez où votre démonstration
utilise l'axiome du choix.
\item[(b)] Soit $ \mathbf{Mat}_k $ la catégorie dont les objets sont les
entiers naturels et dont les morphismes $ m \to n $ sont les matrices $ n \times m $ sur
$ k $, composées par multiplication matricielle. Démontrez que $ \mathbf{Mat}_k $ est
équivalente, mais non isomorphe, à la catégorie des $ k $-espaces vectoriels de dimension finie.
\end{enumerate}

\end{problem}

\noindent L'équivalence, et non l'isomorphisme, est la bonne notion de \og{} même catégorie \fg{} ---
l'une des premières leçons durement acquises du sujet, puisqu'elle impose d'abandonner l'égalité
des objets au profit de l'isomorphisme partout. La partie (b) rend exacte la devise la plus
profonde de l'algèbre linéaire : choisir des bases transforme toute l'algèbre linéaire de
dimension finie en arithmétique matricielle, et la catégorie des matrices est un \textit{squelette} ---
une catégorie équivalente sans objets redondants.

\begin{refs}
\item Contexte : Équivalence de catégories --- Wikipédia (\url{https://fr.wikipedia.org/wiki/%C3%89quivalence_de_cat%C3%A9gories})
\item Lecture : Emily Riehl, \textit{Category Theory in Context}, ch. 1
\item Lecture : Tom Leinster, \textit{Basic Category Theory}, ch. 1
\end{refs}

\bigskip

\begin{problem}[{Tout égaliseur est un monomorphisme --- Licence, short}]
\textbf{CAT-16.} \textbf{Problème.} Un \textit{égaliseur} de $ f, g\colon A \to B $ est un morphisme
$ e\colon E \to A $ avec $ f \circ e = g \circ e $ qui est universel parmi de tels morphismes.
Démontrez que tout égaliseur est un monomorphisme, et que dans $ \mathbf{Set} $ l'égaliseur de
$ f $ et $ g $ est l'inclusion de $ \{\, a \in A : f(a) = g(a) \,\} $.
\end{problem}

\noindent Les égaliseurs sont le mot catégorique pour \og{} ensemble des solutions d'une
équation \fg{}, et la propriété de monomorphisme est le résidu, en termes de flèches, du fait qu'un
ensemble de solutions est un sous-ensemble. L'énoncé dual --- tout coégaliseur est un épimorphisme ---
est offert par la dualité de CAT-04, et les coégaliseurs sont la manière dont les quotients par
une relation entrent dans une catégorie qui n'a aucun élément à identifier.

\begin{refs}
\item Contexte : Égaliseur (mathématiques) --- Wikipédia (\url{https://fr.wikipedia.org/wiki/%C3%89galiseur_%28math%C3%A9matiques%29})
\item Contexte : Coégaliseur --- Wikipédia (en anglais) (\url{https://en.wikipedia.org/wiki/Coequalizer})
\end{refs}

\bigskip

\begin{problem}[{Trancher une catégorie au-dessus d'un objet --- Licence, short}]
\textbf{CAT-17.} \textbf{Problème.} Pour un ensemble $ B $, soit $ \mathbf{Set}/B $ la catégorie
dont les objets sont les applications $ p\colon X \to B $ et dont les morphismes $ p \to q $ sont
les applications $ h $ telles que $ q \circ h = p $. Démontrez que $ \mathbf{Set}/B $ est
équivalente à la catégorie de foncteurs $ \mathbf{Set}^{B} $, où $ B $ est vu comme une catégorie
discrète --- c'est-à-dire à la catégorie des familles d'ensembles indexées par $ B $.
\end{problem}

\noindent Un sens prend les fibres, $ p \mapsto (p^{-1}(b))_{b \in B} $ ; l'autre prend
les réunions disjointes. Cette équivalence explique pourquoi \og{} une famille d'ensembles \fg{} et \og{} un
ensemble s'envoyant sur un ensemble d'indices \fg{} sont employés indifféremment, et sous sa forme
faisceautique --- espaces au-dessus de $ B $ contre faisceaux sur $ B $ --- elle devient la
construction de l'espace étalé.

\begin{refs}
\item Contexte : Catégorie virgule --- Wikipédia (en anglais) (\url{https://en.wikipedia.org/wiki/Comma_category})
\item Contexte : Catégorie de foncteurs --- Wikipédia (\url{https://fr.wikipedia.org/wiki/Cat%C3%A9gorie_de_foncteurs})
\end{refs}

\bigskip

\begin{problem}[{Catégories virgules --- Licence}]
\textbf{CAT-18.} \textbf{Problème.} Étant donné des foncteurs $ F\colon \mathcal{A} \to \mathcal{C} $ et
$ G\colon \mathcal{B} \to \mathcal{C} $, la \textit{catégorie virgule}
$ (F \downarrow G) $ a pour objets les triplets $ (A, B, f) $ avec $ A $ dans
$ \mathcal{A} $, $ B $ dans $ \mathcal{B} $ et $ f\colon F(A) \to G(B) $, et pour
morphismes $ (A, B, f) \to (A', B', f') $ les couples $ (a\colon A \to A',\ b\colon B \to B') $
tels que $ G(b) \circ f = f' \circ F(a) $.
\begin{enumerate}
\item[(a)] Vérifiez que $ (F \downarrow G) $ est une catégorie et que les deux projections vers
$ \mathcal{A} $ et $ \mathcal{B} $ sont des foncteurs.
\item[(b)] Identifiez comme catégories virgules : la tranche $ \mathcal{C}/C $ au-dessus d'un objet
$ C $ ; la cotranche $ C/\mathcal{C} $ ; et la catégorie des flèches, dont les objets sont
tous les morphismes de $ \mathcal{C} $.
\item[(c)] Démontrez que $ \mathrm{id}_C $ est l'objet final de $ \mathcal{C}/C $, et qu'un produit
binaire dans $ \mathcal{C}/C $ est exactement un produit fibré dans $ \mathcal{C} $.
\item[(d)] On appelle $ u\colon C \to G(B_0) $ une \textit{flèche universelle} de $ C $ vers $ G $ si
c'est un objet initial de $ (C \downarrow G) $. Démontrez que $ G $ admet un adjoint à gauche
si et seulement si un tel objet initial existe pour tout $ C $, et confrontez l'énoncé à
l'adjonction du groupe libre de CAT-09.
\item[(e)] Supposez que $ \mathcal{C} $ possède les produits fibrés. Démontrez que pour
$ f\colon C \to D $ le foncteur image réciproque
$ f^{*}\colon \mathcal{C}/D \to \mathcal{C}/C $ est adjoint à droite de la composition avec
$ f $.
\end{enumerate}

\end{problem}

\noindent Les catégories virgules ont été introduites par F. William Lawvere dans sa thèse
de 1963 à Columbia, \textit{Functorial Semantics of Algebraic Theories} ; ce nom étrange garde la
trace de la notation originale $ (F, G) $, une virgule entre deux foncteurs. Ce sont elles qui
transforment une propriété universelle en un banal objet initial, et c'est pourquoi la partie (d)
réduit toute la théorie des adjonctions à une question portant sur un seul objet d'une seule
catégorie --- et pourquoi le théorème du foncteur adjoint (CAT-21) se démontre en chassant les
objets initiaux. La partie (e) est le début de la théorie des types dépendants lue
catégoriquement : la substitution est un produit fibré.

\begin{refs}
\item Contexte : Catégorie virgule --- Wikipédia (en anglais) (\url{https://en.wikipedia.org/wiki/Comma_category})
\item Contexte : Propriété universelle --- Wikipédia (\url{https://fr.wikipedia.org/wiki/Propri%C3%A9t%C3%A9_universelle})
\item Lecture : Saunders Mac Lane, \textit{Categories for the Working Mathematician}, ch. II \S{}6 et ch. IV
\item Lecture : Emily Riehl, \textit{Category Theory in Context}, ch. 4
\end{refs}

\bigskip

\begin{problem}[{Correspondances de Galois : des adjonctions que l'on voit --- Licence}]
\textbf{CAT-26.} \textbf{Problème.} Soient $ (P, \le) $ et $ (Q, \le) $ des ensembles ordonnés, chacun vu
comme une catégorie ayant exactement un morphisme $ p \to p' $ lorsque $ p \le p' $ et aucun
sinon. Une \textit{correspondance de Galois} (monotone) est un couple d'applications croissantes
$ F \colon P \to Q $ et $ G \colon Q \to P $ telles que
\[ F(p) \le q \quad \Longleftrightarrow \quad p \le G(q) \qquad \text{pour tous } p \in P,\ q \in Q. \]
\begin{enumerate}
\item[(a)] Démontrez que c'est exactement une adjonction $ F \dashv G $ entre les catégories
associées, et que l'unité et la counité de CAT-09 deviennent les deux inégalités
$ p \le G(F(p)) $ et $ F(G(q)) \le q $. Démontrez que chacune de $ F, G $ détermine
l'autre de manière unique.
\item[(b)] Démontrez que $ F $ préserve toute borne supérieure existant dans $ P $ et $ G $ toute
borne inférieure existant dans $ Q $ --- l'ombre, sur les ensembles ordonnés, de \og{} les adjoints à
gauche préservent les colimites, les adjoints à droite les limites \fg{} (CAT-10). Montrez par un
exemple que la réciproque est fausse sans hypothèse de complétude, et énoncez le théorème du
foncteur adjoint pour les treillis complets que cette hypothèse achète.
\item[(c)] Démontrez que $ GF $ est un opérateur de fermeture sur $ P $ (extensif, croissant,
idempotent) et $ FG $ un opérateur d'intérieur sur $ Q $, et que $ F $ et $ G $ se
restreignent en des isomorphismes d'ordre mutuellement inverses entre les éléments
$ GF $-fermés de $ P $ et les éléments $ FG $-fermés de $ Q $.
\item[(d)] Identifiez chacune des situations suivantes comme une correspondance de Galois et décrivez
ses points fixes : les sous-groupes de $ \mathrm{Gal}(L/K) $ face aux corps intermédiaires
d'une extension galoisienne ; les parties de l'espace affine de dimension $ n $ face aux idéaux
de $ k[x_1, \dots, x_n] $, via $ V $ et $ I $ ; et, pour une application
$ f \colon X \to Y $, la chaîne $ f_{!} \dashv f^{-1} \dashv f_{*} $ entre ensembles de
parties, où $ f_{!} $ est l'image directe et
$ f_{*}(A) = \{\, y \in Y : f^{-1}(y) \subseteq A \,\} $. Remarquez que les deux premières sont
\textit{décroissantes}, et dites précisément comment les présenter comme des correspondances
monotones ; puis expliquez la lecture que Lawvere fait de la troisième,
$ \exists \dashv \text{substitution} \dashv \forall $.
\end{enumerate}

\end{problem}

\noindent Øystein Ore a nommé et généralisé la notion dans \og{} Galois connexions \fg{} (1944),
en abstrayant la correspondance que Galois avait utilisée en 1830 entre sous-groupes et
sous-corps, et les théoriciens des treillis l'ont ensuite trouvée partout --- opérateurs de
fermeture, analyse formelle de concepts, syntaxe contre sémantique. La théorie des catégories a
absorbé tout le sujet en une ligne : une correspondance de Galois est une adjonction entre
ensembles ordonnés, et tout phénomène propre aux adjonctions y est visible sans la comptabilité
des 2-cellules, puisqu'un ensemble ordonné n'a aucune 2-cellule non triviale à suivre. Lawvere a
ensuite renversé le sens de l'enseignement dans \og{} Adjointness in Foundations \fg{} (1969), en lisant
les quantificateurs de la logique comme des adjoints de la substitution --- et c'est là que
commencent le traitement catégorique de la logique et le langage interne d'un topos (CAT-31).

\begin{refs}
\item Contexte : Correspondance de Galois --- Wikipédia (\url{https://fr.wikipedia.org/wiki/Correspondance_de_Galois})
\item Contexte : Opérateur de fermeture --- Wikipédia (en anglais) (\url{https://en.wikipedia.org/wiki/Closure_operator})
\item Article : O. Ore, \og{} Galois connexions \fg{} (1944), Trans. Amer. Math. Soc. 55
\item Lecture : Tom Leinster, \textit{Basic Category Theory}, ch. 2
\end{refs}

\bigskip

\begin{problem}[{Algèbres initiales : le lemme de Lambek et les entiers naturels --- Licence}]
\textbf{CAT-27.} \textbf{Problème.} Pour un endofoncteur $ F \colon \mathcal{C} \to \mathcal{C} $, une
\textit{$ F $-algèbre} est un couple $ (A, \alpha) $ avec $ \alpha \colon F(A) \to A $, et
un morphisme $ (A, \alpha) \to (B, \beta) $ est une application $ h \colon A \to B $ telle que
$ h \circ \alpha = \beta \circ F(h) $.
\begin{enumerate}
\item[(a)] Vérifiez que les $ F $-algèbres et ces morphismes forment une catégorie, et démontrez le
\textit{lemme de Lambek} : si $ (A, \alpha) $ est une $ F $-algèbre initiale alors $ \alpha $
est un isomorphisme, de sorte que $ A $ est un point fixe de $ F $.
\item[(b)] Prenez $ \mathcal{C} = \mathbf{Set} $ et $ F(X) = 1 + X $, la réunion disjointe avec un
singleton. Démontrez que $ \mathbb{N} $, muni de l'application $ 1 + \mathbb{N} \to \mathbb{N} $
donnée par zéro et le successeur, est une $ F $-algèbre initiale, et démontrez que l'initialité
est exactement le principe de définition par récursion primitive : pour tout ensemble $ B $,
tout élément $ b \in B $ et toute application $ u \colon B \to B $ il existe une unique
$ f \colon \mathbb{N} \to B $ avec $ f(0) = b $ et $ f(n+1) = u(f(n)) $.
\item[(c)] Démontrez que l'algèbre initiale de $ F(X) = 1 + A \times X $ est l'ensemble $ A^{*} $ des
listes finies d'éléments de $ A $, et identifiez les uniques morphismes d'algèbres issus d'elle
au fold familier. Démontrez ensuite que le foncteur parties covariant $ \mathcal{P} $ sur
$ \mathbf{Set} $ n'a aucune algèbre initiale --- et même aucune algèbre dont l'application de
structure soit un isomorphisme.
\item[(d)] Compréhension : énoncez la construction d'Adámek --- lorsque $ \mathcal{C} $ possède un objet
initial $ 0 $ et les colimites de $ \omega $-chaînes, et que $ F $ les préserve, l'algèbre
initiale est la colimite de $ 0 \to F(0) \to F^2(0) \to \cdots $ --- et expliquez pourquoi la
notion duale, une \textit{coalgèbre} finale $ A \to F(A) $, est le bon foyer pour les flux infinis
et les autres objets spécifiés par ce que l'on peut en observer plutôt que par la manière dont on
les construit.
\end{enumerate}

\end{problem}

\noindent Joachim Lambek a démontré le lemme du point fixe dans \og{} A fixpoint theorem for
complete categories \fg{} (1968), en cherchant un foyer catégorique aux définitions récursives. Le
cercle qu'il referme est l'un des plus jolis du sujet : l'objet des entiers naturels de Lawvere,
en 1964, est précisément l'algèbre initiale de $ X \mapsto 1 + X $, si bien que récurrence et
récursion cessent d'être des axiomes sur les nombres pour devenir une propriété universelle qu'une
catégorie possède ou non --- et la partie (c) exhibe un foncteur pour lequel aucun tel objet ne peut
exister, le théorème de Cantor déguisé. Les langages de programmation ont repris l'idée telle
quelle : un type de données algébrique est une algèbre initiale, \texttt{fold} est l'unique
morphisme qui en sort, et les codonnées des flux et des processus racontent l'histoire duale des
coalgèbres finales.

\begin{refs}
\item Contexte : Algèbre initiale --- Wikipédia (en anglais) (\url{https://en.wikipedia.org/wiki/Initial_algebra})
\item Contexte : F-algèbre --- Wikipédia (en anglais) (\url{https://en.wikipedia.org/wiki/F-algebra})
\item Contexte : Objet nombres naturels --- Wikipédia (\url{https://fr.wikipedia.org/wiki/Objet_nombres_naturels})
\item Article : J. Lambek, \og{} A fixpoint theorem for complete categories \fg{} (1968), Mathematische Zeitschrift 103
\end{refs}

\bigskip

\begin{problem}[{Le lemme de recollement pour les produits fibrés --- Licence, short}]
\textbf{CAT-28.} \textbf{Problème.} Soient
\[ \begin{array}{ccccc}
A & \longrightarrow & B & \longrightarrow & C \\
\downarrow & & \downarrow & & \downarrow \\
D & \longrightarrow & E & \longrightarrow & F
\end{array} \]
deux carrés commutatifs dans une catégorie, partageant l'arête verticale du milieu. Démontrez que
si le carré de droite est cartésien, alors le carré de gauche est cartésien si et seulement si le
rectangle extérieur l'est, et donnez un exemple montrant que l'hypothèse sur le carré de droite ne
peut pas être supprimée.
\end{problem}

\noindent C'est le premier véritable outil de la chasse au diagramme : il permet
d'assembler de grands produits fibrés à partir de petits et de reconnaître les petits à
l'intérieur des grands, et on l'emploie constamment en théorie de la descente, en théorie des
fibrations et en sémantique catégorique, où substituer dans une formule revient à prendre un
produit fibré le long d'une application (CAT-18(e)). L'énoncé dual pour les sommes amalgamées est
ce qui fait fonctionner les arguments de recollement à la van Kampen.

\begin{refs}
\item Contexte : Produit fibré --- Wikipédia (\url{https://fr.wikipedia.org/wiki/Produit_fibr%C3%A9})
\item Contexte : Limite (théorie des catégories) --- Wikipédia (\url{https://fr.wikipedia.org/wiki/Limite_%28th%C3%A9orie_des_cat%C3%A9gories%29})
\end{refs}

\bigskip


\section*{Master}

\begin{problem}[{Foncteurs représentables : un échauffement pour Yoneda --- Master}]
\textbf{CAT-07.} \textbf{Problème.} Un foncteur $ F\colon \mathcal{C} \to \mathbf{Set} $ est
\textit{représentable} s'il est naturellement isomorphe à $ \mathrm{Hom}(A, -) $ pour un certain
objet $ A $, l'\textit{objet représentant}.
\begin{enumerate}
\item[(a)] Démontrez que le foncteur d'oubli
$ \mathbf{Grp} \to \mathbf{Set} $ est représenté par $ \mathbb{Z} $, le foncteur d'oubli
$ \mathbf{Ring} \to \mathbf{Set} $ par l'anneau de polynômes $ \mathbb{Z}[x] $, et le
foncteur d'oubli $ \mathbf{Top} \to \mathbf{Set} $ par l'espace à un point.
\item[(b)] Démontrez
que le foncteur parties contravariant $ \mathcal{P}\colon \mathbf{Set}^{\mathrm{op}} \to
\mathbf{Set} $ est représenté par un ensemble à deux éléments.
\item[(c)] Démontrez le théorème de Cayley --- tout
groupe se plonge dans un groupe symétrique --- et expliquez en quoi votre démonstration est le cas
à un objet d'un énoncé général sur les foncteurs $ \mathrm{Hom} $.
\end{enumerate}

\end{problem}

\noindent La représentabilité est le pont entre \og{} une construction \fg{} et \og{} un objet \fg{} :
l'objet représentant porte un élément universel à partir duquel le foncteur tout entier peut être
reconstitué. Grothendieck a bâti la géométrie algébrique moderne sur cette idée --- un schéma
s'étudie à travers son foncteur des points $ \mathrm{Hom}(-, X) $. La partie (c) est un piège
délibérément tendu un siècle trop tôt : le théorème de Cayley de 1854 est le lemme de Yoneda
(CAT-08) vu par le trou d'une serrure à un seul objet.

\begin{refs}
\item Contexte : Foncteur représentable --- Wikipédia (\url{https://fr.wikipedia.org/wiki/Foncteur_repr%C3%A9sentable})
\item Contexte : Théorème de Cayley --- Wikipédia (\url{https://fr.wikipedia.org/wiki/Th%C3%A9or%C3%A8me_de_Cayley})
\item Lecture : Emily Riehl, \textit{Category Theory in Context}, ch. 2
\end{refs}

\bigskip

\begin{problem}[{Le lemme de Yoneda --- Master}]
\textbf{CAT-08.} \textbf{Problème.} Soit $ \mathcal{C} $ une catégorie localement petite.
\begin{enumerate}
\item[(a)] Démontrez le lemme de Yoneda : pour tout foncteur $ F\colon \mathcal{C} \to \mathbf{Set} $ et
tout objet $ A $, l'application envoyant une transformation naturelle
$ \alpha\colon \mathrm{Hom}(A,-) \Rightarrow F $ sur l'élément
$ \alpha_A(\mathrm{id}_A) \in F(A) $ est une bijection
\[ \mathrm{Nat}(\mathrm{Hom}(A,-),\, F) \;\cong\; F(A), \]
naturelle en $ A $ comme en $ F $.
\item[(b)] Déduisez-en que le plongement de Yoneda $ A \mapsto
\mathrm{Hom}(-, A) $ est un foncteur plein et fidèle de $ \mathcal{C} $ vers la
catégorie des foncteurs contravariants à valeurs ensemblistes sur $ \mathcal{C} $.
\item[(c)] Déduisez-en que si
$ \mathrm{Hom}(A,-) \cong \mathrm{Hom}(B,-) $ naturellement, alors $ A \cong B $ --- et
concluez que les objets définis par des propriétés universelles (CAT-03, CAT-04) sont déterminés
à isomorphisme près par le foncteur qu'ils représentent.
\end{enumerate}

\end{problem}

\noindent Le lemme est né d'une conversation parisienne de 1954 entre Nobuo Yoneda et
Saunders Mac Lane, qui y attacha plus tard le nom de Yoneda ; il apparaît pleinement formé dans
les séminaires de Grothendieck peu après. Son contenu est une philosophie accompagnée d'une
démonstration : un objet est exactement le système de toutes les applications qui vont vers lui ---
la version mathématique du fait de connaître une chose par ses relations. Presque tout ce qui
suit sur cette page (adjonctions, limites, monades) est le lemme de Yoneda en tenue de travail.

\begin{refs}
\item Contexte : Lemme de Yoneda --- Wikipédia (\url{https://fr.wikipedia.org/wiki/Lemme_de_Yoneda})
\item Lecture : Emily Riehl, \textit{Category Theory in Context}, ch. 2
\item Lecture : Saunders Mac Lane, \textit{Categories for the Working Mathematician}, ch. III
\end{refs}

\bigskip

\begin{problem}[{Adjonctions : libre $\dashv$ oubli --- Master}]
\textbf{CAT-09.} \textbf{Problème.} Des foncteurs $ F\colon \mathcal{C} \to \mathcal{D} $ et
$ G\colon \mathcal{D} \to \mathcal{C} $ forment une adjonction $ F \dashv G $ s'il existe une
bijection $ \mathrm{Hom}_{\mathcal{D}}(F A, B) \cong \mathrm{Hom}_{\mathcal{C}}(A, G B) $,
naturelle en $ A $ et en $ B $.
\begin{enumerate}
\item[(a)] Démontrez que cette définition équivaut à la donnée de
transformations naturelles $ \eta\colon \mathrm{id} \Rightarrow GF $ (l'unité) et
$ \varepsilon\colon FG \Rightarrow \mathrm{id} $ (la counité) satisfaisant les deux identités
triangulaires.
\item[(b)] Vérifiez les adjonctions : groupe libre $ \dashv $ oubli
$ \mathbf{Grp} \to \mathbf{Set} $ ; $ k $-espace vectoriel libre $ \dashv $ oubli ; et la
chaîne discret $ \dashv $ oubli $ \dashv $ grossier pour les espaces topologiques.
\item[(c)] À l'aide du lemme de Yoneda, démontrez qu'un adjoint à gauche d'un foncteur donné, lorsqu'il
existe, est unique à isomorphisme naturel près.
\end{enumerate}

\end{problem}

\noindent Daniel Kan a isolé les foncteurs adjoints en 1958, en abstrayant la relation
tenseur--hom de l'algèbre et son propre travail en homotopie simpliciale ; en moins de dix ans Mac
Lane pouvait écrire la devise de son manuel : \og{} Adjoint functors arise everywhere \fg{} (les foncteurs
adjoints surgissent partout). Les couples libre/oubli de (b) en sont le modèle --- toute
construction \og{} libre \fg{} en algèbre est un adjoint à gauche, et sa propriété universelle est l'unité
de l'adjonction.

\begin{refs}
\item Contexte : Foncteur adjoint --- Wikipédia (\url{https://fr.wikipedia.org/wiki/Foncteur_adjoint})
\item Article : D. M. Kan, \og{} Adjoint functors \fg{} (1958), Trans. Amer. Math. Soc. 87
\item Lecture : Saunders Mac Lane, \textit{Categories for the Working Mathematician}, ch. IV
\end{refs}

\bigskip

\begin{problem}[{Limites, colimites, et ce que préservent les adjoints --- Master}]
\textbf{CAT-10.} \textbf{Problème.} Une \textit{limite} d'un diagramme $ D\colon \mathcal{J} \to
\mathcal{C} $ (avec $ \mathcal{J} $ petite) est un objet final parmi les cônes au-dessus de
$ D $ ; une \textit{colimite} est la notion duale.
\begin{enumerate}
\item[(a)] Montrez que les produits, les égaliseurs, les produits fibrés et les objets finaux sont des
limites sur des diagrammes de formes convenables, et démontrez qu'une catégorie possédant tous les
petits produits et tous les égaliseurs possède toutes les petites limites.
\item[(b)] Démontrez que pour tout objet $ A $ le
foncteur $ \mathrm{Hom}(A, -)\colon \mathcal{C} \to \mathbf{Set} $ préserve toutes les limites.
\item[(c)] Démontrez que les adjoints à droite préservent les limites et, dualement, que les adjoints à
gauche préservent les colimites.
\item[(d)] Appliquez (c) par contraposition : le foncteur d'oubli $ U\colon \mathbf{Grp} \to
\mathbf{Set} $ envoie le coproduit $ \mathbb{Z} * \mathbb{Z} $ (CAT-03) sur un ensemble bien
plus grand que la réunion disjointe de deux copies de $ \mathbb{Z} $ ; concluez que $ U $ n'a
aucun adjoint à droite.
\end{enumerate}

\end{problem}

\noindent Les limites et colimites, en toute généralité, sont dues au même cercle --- encore
l'article de Kan de 1958 --- et la partie (c), \og{} RAPL \fg{} en anglais, est le cheval de trait quotidien
du sujet : elle explique en une ligne pourquoi les constructions libres malmènent les produits,
pourquoi les ensembles sous-jacents des produits sont des produits, et (sous sa forme raffinée,
les théorèmes du foncteur adjoint de Freyd) quand un foncteur préservant les limites doit avoir un
adjoint. La partie (b) est le point de vue de Yoneda sur les limites : elles se calculent point par
point, sous le regard de chaque objet.

\begin{refs}
\item Contexte : Limite (théorie des catégories) --- Wikipédia (\url{https://fr.wikipedia.org/wiki/Limite_%28th%C3%A9orie_des_cat%C3%A9gories%29})
\item Lecture : Emily Riehl, \textit{Category Theory in Context}, ch. 3--4
\item Lecture : Saunders Mac Lane, \textit{Categories for the Working Mathematician}, ch. V
\end{refs}

\bigskip

\begin{problem}[{Des monades à partir des adjonctions --- Master}]
\textbf{CAT-11.} \textbf{Problème.} Une \textit{monade} sur $ \mathcal{C} $ est un foncteur
$ T\colon \mathcal{C} \to \mathcal{C} $ muni de transformations naturelles
$ \eta\colon \mathrm{id} \Rightarrow T $ et $ \mu\colon T^2 \Rightarrow T $ satisfaisant
l'associativité $ \mu \circ T\mu = \mu \circ \mu T $ et les lois d'unité
$ \mu \circ T\eta = \mu \circ \eta T = \mathrm{id} $.
\begin{enumerate}
\item[(a)] Étant donné une adjonction
$ F \dashv G $ d'unité $ \eta $ et de counité $ \varepsilon $, démontrez que
$ T = GF $, avec l'unité $ \eta $ et la multiplication $ \mu = G \varepsilon F $, est une
monade.
\item[(b)] Identifiez la monade sur $ \mathbf{Set} $ induite par l'adjonction du monoïde libre ---
montrez que $ T(X) = X^{*} $, l'ensemble des mots finis sur $ X $ --- et vérifiez directement
que le foncteur parties, avec l'unité $ x \mapsto \{x\} $ et la multiplication donnée par la
réunion, est une monade.
\item[(c)] Démontrez la réciproque de (a) : toute monade provient d'une adjonction, en construisant
la catégorie de Kleisli $ \mathcal{C}_T $ et en exhibant l'adjonction qui induit
$ T $.
\end{enumerate}

\end{problem}

\noindent Les monades ont été distillées par Godement (1958, sous le nom de
\og{} constructions standard \fg{}) pour les résolutions faisceautiques ; Kleisli, et Eilenberg avec
Moore, ont répondu en 1965 à la question \og{} toute monade ? \fg{} par deux adjonctions canoniques
différentes. Trente ans plus tard, Moggi remarquait que les effets computationnels --- exceptions,
état, non-déterminisme, entrées/sorties --- sont des monades, et le concept passait dans la pratique
de la programmation ; la monade des listes de (b) tient en une ligne de définition en Haskell. Peu
d'abstractions ont voyagé aussi loin.

\begin{refs}
\item Contexte : Monade (théorie des catégories) --- Wikipédia (\url{https://fr.wikipedia.org/wiki/Monade_%28th%C3%A9orie_des_cat%C3%A9gories%29})
\item Contexte : Monade (informatique) --- Wikipédia (\url{https://fr.wikipedia.org/wiki/Monade_%28informatique%29})
\item Lecture : Emily Riehl, \textit{Category Theory in Context}, ch. 5
\item Article : E. Moggi, \og{} Notions of computation and monads \fg{} (1991), Information and Computation 93
\end{refs}

\bigskip

\begin{problem}[{Catégories cartésiennes fermées et $\lambda$-calcul --- Master}]
\textbf{CAT-12.} \textbf{Problème.} Une catégorie à produits finis est \textit{cartésienne
fermée} (une CCC) si pour tous objets $ A, B $ il existe un objet exponentiel $ B^A $
muni d'une évaluation $ \mathrm{ev}\colon B^A \times A \to B $ telle que tout
$ f\colon C \times A \to B $ se factorise en $ f = \mathrm{ev} \circ (\lambda f \times
\mathrm{id}_A) $ pour un unique $ \lambda f\colon C \to B^A $.
\begin{enumerate}
\item[(a)] Démontrez que
$ \mathbf{Set} $ est cartésienne fermée, et que la propriété universelle revient à la bijection
de curryfication $ \mathrm{Hom}(C \times A,\, B) \cong \mathrm{Hom}(C,\, B^A) $.
\item[(b)] Démontrez que $ \mathbf{Grp} $ n'est pas cartésienne fermée.
\item[(c)] Le $\lambda$-calcul simplement typé
a des types construits à partir de types de base par $ \times $ et $ \to $, et des termes
construits par appariement, projection, application et $\lambda$-abstraction. Vérifiez soigneusement que
les types-comme-objets et les termes-modulo-$ \beta\eta $-équivalence-comme-morphismes
fournissent une catégorie cartésienne fermée, la $\lambda$-abstraction réalisant la propriété universelle
de (a).
\end{enumerate}

\end{problem}

\noindent La partie (c) est un tiers de la correspondance de Curry--Howard--Lambek : les
démonstrations en logique intuitionniste, les programmes en $\lambda$-calcul typé et les morphismes des
catégories cartésiennes fermées sont trois notations pour les mêmes objets. Lambek a établi le
pied catégorique vers 1980 ; Lawvere s'était déjà servi des CCC pour reformuler la logique dans
les années 1960. Ce triangle est l'ossature théorique des langages fonctionnels typés et des
assistants de preuve qui apparaissent dans la bande Recherche.

\begin{refs}
\item Contexte : Catégorie cartésienne fermée --- Wikipédia (\url{https://fr.wikipedia.org/wiki/Cat%C3%A9gorie_cart%C3%A9sienne_ferm%C3%A9e})
\item Contexte : Correspondance de Curry--Howard --- Wikipédia (\url{https://fr.wikipedia.org/wiki/Correspondance_de_Curry-Howard})
\item Contexte : Lambda-calcul simplement typé --- Wikipédia (\url{https://fr.wikipedia.org/wiki/Lambda-calcul_simplement_typ%C3%A9})
\item Article : J. Lambek, \og{} From $\lambda$-calculus to cartesian closed categories \fg{} (1980), dans \textit{To H. B. Curry: Essays on Combinatory Logic, Lambda Calculus and Formalism}
\end{refs}

\bigskip

\begin{problem}[{Les limites dans une catégorie de foncteurs se calculent point par point --- Master, short}]
\textbf{CAT-19.} \textbf{Problème.} Soit $ \mathcal{C} $ petite et soit $ \mathcal{D} $
possédant toutes les petites limites. Démontrez que la catégorie de foncteurs
$ [\mathcal{C}, \mathcal{D}] $ possède toutes les petites limites, et qu'elles se calculent
objet par objet : pour un diagramme $ F\colon \mathcal{J} \to [\mathcal{C}, \mathcal{D}] $ et
tout objet $ c $ de $ \mathcal{C} $,
\[ \left( \lim_{\mathcal{J}} F \right)(c) \;\cong\; \lim_{\mathcal{J}} \big( F(-)(c) \big). \]
\end{problem}

\noindent C'est le cheval de trait derrière tout argument du type \og{} calculons objet par
objet \fg{} : noyaux de transformations naturelles, limites de préfaisceaux, limites de diagrammes de
modules. L'énoncé dual vaut pour les colimites, et l'asymétrie qui rend le sujet intéressant
apparaît un étage plus haut --- les extensions de Kan \textit{à gauche} et les réalisations
géométriques ne se calculent pas point par point en ce sens naïf.

\begin{refs}
\item Contexte : Catégorie de foncteurs --- Wikipédia (\url{https://fr.wikipedia.org/wiki/Cat%C3%A9gorie_de_foncteurs})
\item Contexte : Limite (théorie des catégories) --- Wikipédia (\url{https://fr.wikipedia.org/wiki/Limite_%28th%C3%A9orie_des_cat%C3%A9gories%29})
\end{refs}

\bigskip

\begin{problem}[{Tout préfaisceau est une colimite de représentables --- Master, short}]
\textbf{CAT-20.} \textbf{Problème.} Soit $ \mathcal{C} $ une petite catégorie, soit
$ X\colon \mathcal{C}^{\mathrm{op}} \to \mathbf{Set} $ un préfaisceau, et soit $ \int X $ sa
\textit{catégorie des éléments}, dont les objets sont les couples $ (c, x) $ avec $ x \in X(c) $
et dont les morphismes $ (c, x) \to (c', x') $ sont les applications $ u\colon c \to c' $ avec
$ X(u)(x') = x $. Démontrez que $ X $ est la colimite du diagramme
$ \int X \to [\mathcal{C}^{\mathrm{op}}, \mathbf{Set}] $ envoyant $ (c, x) $ sur le
préfaisceau représentable $ \mathrm{Hom}(-, c) $.
\end{problem}

\noindent Le théorème de densité dit que le plongement de Yoneda est aussi généreux que
possible : les représentables engendrent tout par colimites, si bien qu'un foncteur préservant les
colimites et défini sur les préfaisceaux est déterminé par ses valeurs sur $ \mathcal{C} $.
C'est le moteur de la devise de la \og{} cocomplétion libre \fg{} et de la recette standard pour construire
des réalisations géométriques --- les ensembles simpliciaux sont des colimites de simplexes, et tout
choix de \og{} ce que signifie un simplexe \fg{} s'étend de manière unique.

\begin{refs}
\item Contexte : Préfaisceau (théorie des catégories) --- Wikipédia (\url{https://fr.wikipedia.org/wiki/Pr%C3%A9faisceau_%28th%C3%A9orie_des_cat%C3%A9gories%29})
\item Contexte : Catégorie des éléments --- Wikipédia (en anglais) (\url{https://en.wikipedia.org/wiki/Category_of_elements})
\end{refs}

\bigskip

\begin{problem}[{Le théorème du foncteur adjoint --- Master}]
\textbf{CAT-21.} \textbf{Problème.} D'après CAT-10, un foncteur admettant un adjoint à gauche préserve les
limites. Les théorèmes du foncteur adjoint demandent quand la réciproque est vraie. On dit que
$ G\colon \mathcal{D} \to \mathcal{C} $ satisfait la \textit{condition d'ensemble solution} en un
objet $ C $ s'il existe un \textbf{ensemble} de morphismes $ f_i\colon C \to G(D_i) $ par
lesquels tout morphisme $ C \to G(D) $ se factorise sous la forme $ G(h) \circ f_i $ pour un
certain $ i $ et un certain $ h\colon D_i \to D $.
\begin{enumerate}
\item[(a)] Démontrez le théorème général du foncteur adjoint : si $ \mathcal{D} $ est petitement
complète et localement petite, et si $ G $ préserve les petites limites et satisfait la
condition d'ensemble solution en tout objet, alors $ G $ admet un adjoint à gauche. (Suivant
CAT-18(d), il suffit de produire un objet initial dans chaque catégorie virgule
$ (C \downarrow G) $ ; construisez-le comme une limite convenable et servez-vous d'un argument
d'\og{} absence d'endomorphisme propre \fg{}.)
\item[(b)] Vérifiez les hypothèses pour le foncteur d'oubli $ \mathbf{Grp} \to \mathbf{Set} $ en
exhibant un ensemble solution explicite, et retrouvez l'existence des groupes libres sans
construire un seul mot.
\item[(c)] Montrez que la préservation des limites ne suffit pas : le foncteur d'oubli de la catégorie
des algèbres de Boole complètes et de leurs homomorphismes complets vers $ \mathbf{Set} $
préserve toutes les petites limites, et pourtant n'a aucun adjoint à gauche, parce qu'il n'existe
pas d'algèbre de Boole complète libre sur une infinité dénombrable de générateurs. Localisez
précisément où la condition d'ensemble solution échoue.
\item[(d)] Énoncez le théorème spécial du foncteur adjoint (ajoutez le caractère bien puissant et un
petit ensemble cogénérateur, supprimez la condition d'ensemble solution) et servez-vous-en pour
démontrer que l'inclusion des espaces compacts séparés dans les espaces topologiques admet un
adjoint à gauche --- la compactification de Stone--Čech.
\end{enumerate}

\end{problem}

\noindent Peter Freyd a démontré ces théorèmes au début des années 1960 et les a publiés
dans \textit{Abelian Categories} (1964) ; ils sont la réponse professionnelle à la question \og{} cette
construction existe-t-elle ? \fg{}, remplaçant les constructions au cas par cas par une vérification
sur les limites accompagnée d'une condition de taille. Le contre-exemple de (c) est celui qui
montre que la condition de taille n'a rien de bureaucratique : Hales et Gaifman ont montré
indépendamment en 1964 que les algèbres de Boole complètes libres sur une infinité de générateurs
n'existent pas, de sorte que les hypothèses du théorème sont la frontière honnête de la méthode.
La partie (d) explique pourquoi Stone--Čech est un adjoint à gauche et préserve donc les coproduits
qu'elle est par ailleurs célèbre pour malmener en tout autre sens.

\begin{refs}
\item Contexte : Foncteur adjoint --- Wikipédia (\url{https://fr.wikipedia.org/wiki/Foncteur_adjoint})
\item Contexte : Algèbre de Boole complète --- Wikipédia (en anglais) (\url{https://en.wikipedia.org/wiki/Complete_Boolean_algebra})
\item Contexte : Compactification de Stone--Čech --- Wikipédia (\url{https://fr.wikipedia.org/wiki/Compactification_de_Stone-%C4%8Cech})
\item Lecture : Saunders Mac Lane, \textit{Categories for the Working Mathematician}, ch. V \S{}6
\item Lecture : Emily Riehl, \textit{Category Theory in Context}, ch. 4
\end{refs}

\bigskip

\begin{problem}[{Catégories monoïdales, cohérence et diagrammes de cordes --- Master}]
\textbf{CAT-22.} \textbf{Problème.} Une \textit{catégorie monoïdale} est une catégorie $ \mathcal{V} $ munie d'un
foncteur $ \otimes\colon \mathcal{V} \times \mathcal{V} \to \mathcal{V} $, d'un objet unité
$ I $, et d'isomorphismes naturels
$ \alpha_{A,B,C}\colon (A \otimes B) \otimes C \to A \otimes (B \otimes C) $,
$ \lambda_A\colon I \otimes A \to A $, $ \rho_A\colon A \otimes I \to A $ soumis aux axiomes
du pentagone et du triangle de Mac Lane.
\begin{enumerate}
\item[(a)] Vérifiez que $ (\mathbf{Set}, \times, \{*\}) $, $ (\mathrm{Vect}_k, \otimes, k) $, et
les endofoncteurs d'une catégorie fixée munis de la composition sont monoïdales, la dernière
strictement.
\item[(b)] Un \textit{monoïde} dans $ \mathcal{V} $ est un objet $ M $ muni de
$ \mu\colon M \otimes M \to M $ et $ \eta\colon I \to M $ satisfaisant l'associativité et les
lois d'unité. Démontrez qu'un monoïde dans $ (\mathbf{Set}, \times) $ est un monoïde ordinaire,
qu'un monoïde dans
$ (\mathbf{Ab}, \otimes_{\mathbb{Z}}, \mathbb{Z}) $ est un anneau, et qu'un monoïde dans la
catégorie des endofoncteurs est exactement une monade (CAT-11).
\item[(c)] Démontrez à partir des axiomes que $ \lambda_I = \rho_I\colon I \otimes I \to I $.
\item[(d)] Problème de compréhension : lisez le théorème de cohérence de Mac Lane sous la forme \og{} toute
catégorie monoïdale est monoïdalement équivalente à une catégorie stricte \fg{}, apprenez le calcul
des diagrammes de cordes de Joyal et Street, et servez-vous-en pour démontrer qu'un objet dual ---
un objet $ A^{*} $ muni d'une unité $ I \to A \otimes A^{*} $ et d'une counité
$ A^{*} \otimes A \to I $ satisfaisant les deux identités en zigzag --- est précisément une
adjonction dans la 2-catégorie à un objet déterminée par $ \mathcal{V} $.
\end{enumerate}

\end{problem}

\noindent Mac Lane a écrit le pentagone en 1963 (Bénabou indépendamment), en se demandant
ce qu'il fallait supposer pour que \og{} tous les diagrammes raisonnables commutent \fg{} --- le premier
théorème de cohérence, et le modèle de tous ceux qui ont suivi en théorie des catégories
supérieures. Les diagrammes de cordes transforment ces diagrammes commutatifs en topologie : les
morphismes deviennent des boîtes sur des fils, les axiomes deviennent la liberté de faire glisser
les boîtes, et une démonstration devient une isotopie. Penrose a inventé la notation pour le
calcul tensoriel vers 1971 et Joyal et Street en ont fait un théorème en 1991 ; c'est aujourd'hui
la notation de travail de l'algèbre quantique, de la théorie topologique des champs (CAT-25) et du
ZX-calcul de l'informatique quantique.

\begin{refs}
\item Contexte : Catégorie monoïdale --- Wikipédia (\url{https://fr.wikipedia.org/wiki/Cat%C3%A9gorie_mono%C3%AFdale})
\item Contexte : Théorème de cohérence --- Wikipédia (en anglais) (\url{https://en.wikipedia.org/wiki/Coherence_theorem})
\item Contexte : Diagramme de cordes --- Wikipédia (en anglais) (\url{https://en.wikipedia.org/wiki/String_diagram})
\item Contexte : Objet dual --- Wikipédia (en anglais) (\url{https://en.wikipedia.org/wiki/Dual_object})
\item Article : S. Mac Lane, \og{} Natural associativity and commutativity \fg{} (1963), Rice University Studies 49
\item Article : A. Joyal et R. Street, \og{} The geometry of tensor calculus, I \fg{} (1991), Advances in Mathematics 88
\end{refs}

\bigskip

\begin{problem}[{Extensions de Kan --- Master}]
\textbf{CAT-29.} \textbf{Problème.} Soient $ K \colon \mathcal{A} \to \mathcal{B} $ et
$ F \colon \mathcal{A} \to \mathcal{C} $ des foncteurs. Une \textit{extension de Kan à gauche} de
$ F $ le long de $ K $ est un foncteur $ \mathrm{Lan}_K F \colon \mathcal{B} \to \mathcal{C} $
accompagné d'une transformation naturelle $ \eta \colon F \Rightarrow (\mathrm{Lan}_K F) \circ K $
universelle parmi de tels couples ; l'extension de Kan à droite $ \mathrm{Ran}_K F $ est la
notion duale.
\begin{enumerate}
\item[(a)] Démontrez que si les extensions de Kan à gauche le long de $ K $ existent pour tout $ F $,
elles s'assemblent en un foncteur adjoint à gauche de la restriction
$ K^{*} \colon [\mathcal{B}, \mathcal{C}] \to [\mathcal{A}, \mathcal{C}] $, et démontrez la
réciproque.
\item[(b)] Démontrez la formule ponctuelle : si $ \mathcal{A} $ est petite et $ \mathcal{C} $ est
cocomplète, alors
\[ (\mathrm{Lan}_K F)(b) \;\cong\; \operatorname*{colim}\, \bigl( (K \downarrow b) \to
\mathcal{A} \xrightarrow{\ F\ } \mathcal{C} \bigr), \]
la colimite sur la catégorie virgule de CAT-18 dont les objets sont les couples
$ (a, K(a) \to b) $, et identifiez le $ \eta $ universel.
\item[(c)] Récoltez la devise \og{} tous les concepts sont des extensions de Kan \fg{}. Montrez qu'une colimite de
$ F \colon \mathcal{A} \to \mathcal{C} $ est $ \mathrm{Lan}_K F $ pour $ K $ l'unique
foncteur vers la catégorie finale ; que pour une adjonction $ F \dashv G $ on a
$ G \cong \mathrm{Ran}_F(\mathrm{id}) $, et déterminez en quel sens cette extension de Kan est
\textit{absolue} --- préservée par absolument tout foncteur ; et que pour $ \mathcal{E} $
cocomplète, l'extension de Kan à gauche le long du plongement de Yoneda envoie un foncteur
$ \mathcal{C} \to \mathcal{E} $ sur l'unique foncteur préservant les colimites
$ [\mathcal{C}^{\mathrm{op}}, \mathbf{Set}] \to \mathcal{E} $ qui le prolonge, ce qui est
l'énoncé de cocomplétion libre derrière CAT-20.
\item[(d)] Calculez l'exemple qui a rendu la notion indispensable. Soit $ \Delta $ la catégorie des
simplexes et soit $ r \colon \Delta \to \mathbf{Top} $ envoyant $ [n] $ sur le
$ n $-simplexe topologique standard. Démontrez que l'extension de Kan à gauche de $ r $ le long
du plongement de Yoneda $ \Delta \to \mathbf{sSet} $ est la réalisation géométrique, qu'elle est
adjointe à gauche du foncteur complexe singulier, et que c'est l'adjonction nerf--réalisation
générale produite par (c).
\end{enumerate}

\end{problem}

\noindent Daniel Kan a introduit ces extensions dans l'article même de 1958 qui a nommé les
foncteurs adjoints, parce qu'il en avait besoin pour construire des objets en théorie de
l'homotopie simpliciale. Mac Lane a consacré le dernier chapitre de \textit{Categories for the Working
Mathematician} à ces extensions et a donné à sa section finale le titre délibérément provocateur
\og{} All concepts are Kan extensions \fg{} --- ce qui est le contenu de la partie (c), où limites, adjoints
et cocomplétion libre découlent tous d'une seule propriété universelle. L'adjonction
nerf--réalisation de la partie (d) est la charnière entre combinatoire et topologie, et sous sa
forme moderne c'est ainsi que les $\infty$-catégories de la bande Recherche sont définies tout court.

\begin{refs}
\item Contexte : Extension de Kan --- Wikipédia (\url{https://fr.wikipedia.org/wiki/Extension_de_Kan})
\item Contexte : Ensemble simplicial --- Wikipédia (\url{https://fr.wikipedia.org/wiki/Ensemble_simplicial})
\item Lecture : Saunders Mac Lane, \textit{Categories for the Working Mathematician}, ch. X
\item Lecture : Emily Riehl, \textit{Category Theory in Context}, ch. 6
\end{refs}

\bigskip

\begin{problem}[{Algèbres sur une monade et théorème de Beck --- Master}]
\textbf{CAT-30.} \textbf{Problème.} Soit $ (T, \eta, \mu) $ une monade sur $ \mathcal{C} $ (CAT-11). Une
\textit{algèbre d'Eilenberg--Moore} est un objet $ A $ muni d'une application $ a \colon T(A) \to A $
satisfaisant $ a \circ \eta_A = \mathrm{id}_A $ et $ a \circ \mu_A = a \circ T(a) $ ;
avec les morphismes évidents, elles forment une catégorie $ \mathcal{C}^{T} $.
\begin{enumerate}
\item[(a)] Démontrez que le foncteur d'oubli $ U^{T} \colon \mathcal{C}^{T} \to \mathcal{C} $ admet
un adjoint à gauche dont la monade induite est de nouveau $ T $ --- une deuxième réponse
canonique, à côté de la construction de Kleisli de CAT-11(c), à la question de savoir quelles
adjonctions produisent une monade donnée. Démontrez que $ \mathcal{C}^{T} $ est \textit{finale}
parmi toutes les adjonctions induisant $ T $, tandis que la catégorie de Kleisli est initiale.
\item[(b)] Calculez. Montrez que les algèbres de la monade du monoïde libre sur $ \mathbf{Set} $ sont
exactement les monoïdes et leurs homomorphismes, celles de la monade du groupe libre exactement
les groupes, et celles de la monade des parties exactement les demi-treillis complets pour la
borne supérieure, avec les applications qui la préservent.
\item[(c)] On appelle une adjonction $ F \dashv G $ avec $ G \colon \mathcal{D} \to \mathcal{C} $
\textit{monadique} si le foncteur de comparaison $ \mathcal{D} \to \mathcal{C}^{GF} $ est une
équivalence. Énoncez le théorème de monadicité de Beck --- $ G $ est monadique si et seulement
s'il admet un adjoint à gauche, réfléchit les isomorphismes, et que $ \mathcal{D} $ possède et
$ G $ préserve les coégaliseurs des couples $ G $-scindés --- et démontrez la direction facile,
à savoir qu'un foncteur monadique possède ces trois propriétés.
\item[(d)] Utilisez le théorème deux fois. Démontrez que le foncteur d'oubli des espaces compacts séparés
vers $ \mathbf{Set} $ est monadique, sa monade étant la monade des ultrafiltres, de sorte que
\og{} compact séparé \fg{} est une théorie algébrique dont l'unique opération infinitaire associe à chaque
ultrafiltre de points sa limite. Montrez ensuite que le foncteur d'oubli
$ \mathbf{Top} \to \mathbf{Set} $ n'est \textit{pas} monadique, en exhibant une bijection continue
qui n'est pas un homéomorphisme et en identifiant exactement l'hypothèse du théorème de Beck
qu'elle détruit.
\end{enumerate}

\end{problem}

\noindent En 1965 Eilenberg et Moore, et indépendamment Kleisli, ont montré que toute
monade provient d'une adjonction, de deux manières extrémales ; la thèse de 1967 de Jon Beck à
Columbia a ensuite répondu à la question plus fine : quand une catégorie de structures au-dessus
d'une autre n'est-elle \textit{rien d'autre} que les algèbres d'une monade ? Ce qui rend la réponse
remarquable, c'est ce qu'elle ne mentionne pas : ni signature, ni équations, seulement une
condition sur certains coégaliseurs. La partie (d) est le dividende le plus célèbre du théorème,
la découverte par Ernest Manes que compacité plus séparation est une théorie algébrique dont les
opérations envoient les ultrafiltres de points sur des points --- ce qu'aucune dose de topologie
générale ne laisse deviner, et le début de la vision catégorique de la topologie comme algèbre
au-dessus de $ \mathbf{Set} $.

\begin{refs}
\item Contexte : Théorème de monadicité de Beck --- Wikipédia (en anglais) (\url{https://en.wikipedia.org/wiki/Beck%27s_monadicity_theorem})
\item Contexte : Monade (théorie des catégories) --- Wikipédia (\url{https://fr.wikipedia.org/wiki/Monade_%28th%C3%A9orie_des_cat%C3%A9gories%29})
\item Article : S. Eilenberg et J. C. Moore, \og{} Adjoint functors and triples \fg{} (1965), Illinois J. Math. 9
\item Lecture : Saunders Mac Lane, \textit{Categories for the Working Mathematician}, ch. VI
\end{refs}

\bigskip

\begin{problem}[{Topos, classificateurs de sous-objets et condition de faisceau --- Master}]
\textbf{CAT-31.} \textbf{Problème.} Un \textit{topos élémentaire} est une catégorie possédant les limites finies, les
exponentielles (CAT-12) et un \textit{classificateur de sous-objets} : un monomorphisme
$ \top \colon 1 \to \Omega $ tel que pour tout monomorphisme $ m \colon S \to X $ il existe un
unique $ \chi_m \colon X \to \Omega $ faisant du carré
\[ \begin{array}{ccc}
S & \longrightarrow & 1 \\
\downarrow & & \downarrow \\
X & \longrightarrow & \Omega
\end{array} \]
--- d'arête gauche $ m $, d'arête droite $ \top $ et d'arête inférieure $ \chi_m $ --- un carré
cartésien.
\begin{enumerate}
\item[(a)] Démontrez qu'un ensemble à deux éléments muni de l'application désignant l'un de ses éléments
est un classificateur de sous-objets pour $ \mathbf{Set} $, et qu'un classificateur de
sous-objets, lorsqu'il existe, est unique à isomorphisme canonique près. Démontrez que dans tout
topos $ X \mapsto \mathrm{Sub}(X) $ est un foncteur
$ \mathcal{E}^{\mathrm{op}} \to \mathbf{Set} $ représenté par $ \Omega $.
\item[(b)] Démontrez que les préfaisceaux forment un topos. Pour $ \mathcal{C} $ petite, appelez
\textit{crible} sur $ c $ un ensemble $ S $ de morphismes de but $ c $ stable par
précomposition ; montrez que $ \Omega(c) = \{ \text{cribles sur } c \} $, la restriction le long
de $ u \colon c' \to c $ étant donnée par l'image réciproque des cribles, est un classificateur
de sous-objets pour $ [\mathcal{C}^{\mathrm{op}}, \mathbf{Set}] $.
\item[(c)] Démontrez que $ \mathrm{Sub}(X) $ est une algèbre de Heyting : elle possède les bornes
inférieures et supérieures finies et une opération $ \Rightarrow $ telle que
$ U \wedge V \le W $ si et seulement si $ U \le (V \Rightarrow W) $ --- une adjonction au sens
de CAT-26 --- mais elle ne satisfait pas nécessairement $ U \vee \neg U = 1 $. Exhibez un topos de
préfaisceaux dans lequel le tiers exclu est en défaut, et concluez que la logique interne d'un
topos est intuitionniste.
\item[(d)] Compréhension : un préfaisceau $ F $ sur un espace topologique est un \textit{faisceau}
lorsque, pour tout ouvert $ U $ et tout recouvrement ouvert $ \{ U_i \} $ de $ U $,
l'ensemble $ F(U) $ est l'égaliseur des deux applications évidentes
$ \prod_i F(U_i) \rightrightarrows \prod_{i,j} F(U_i \cap U_j) $. Démontrez que cette unique
condition de limite est exactement la séparation jointe au recollement, démontrez que les
faisceaux forment une sous-catégorie pleine des préfaisceaux stable par limites, et énoncez le
théorème de Lawvere et Tierney selon lequel les faisceaux pour une topologie sur un topos
quelconque forment encore un topos.
\end{enumerate}

\end{problem}

\noindent L'école de Grothendieck a introduit les topos vers 1963, dans SGA 4, comme des
espaces généralisés : des catégories de faisceaux dans lesquelles on pouvait faire de la
cohomologie pour la topologie étale, là où l'\og{} espace \fg{} sous-jacent a trop peu d'ouverts pour être
utile. Puis, en 1970, Lawvere et Tierney ont trouvé les axiomes de la partie (a) --- limites finies,
exponentielles et un unique objet $ \Omega $ représentant les sous-objets --- une description au
premier ordre de ce qu'est une catégorie de faisceaux, sans site ni espace en vue. La conséquence
de la partie (c) est ce qui a changé la logique : tout topos porte un langage interne dont la
logique est intuitionniste, si bien qu'un topos est à la fois un espace généralisé et un univers
mathématique, et que la démonstration par forcing de l'indépendance de l'hypothèse du continu,
due à Cohen, réapparaît comme la construction d'un topos convenable. La page
Fondements (\url{pure-foundations.html}) poursuit du côté logique.

\begin{refs}
\item Contexte : Topos (mathématiques) --- Wikipédia (\url{https://fr.wikipedia.org/wiki/Topos_%28math%C3%A9matiques%29})
\item Contexte : Classificateur de sous-objets --- Wikipédia (en anglais) (\url{https://en.wikipedia.org/wiki/Subobject_classifier})
\item Contexte : Algèbre de Heyting --- Wikipédia (\url{https://fr.wikipedia.org/wiki/Alg%C3%A8bre_de_Heyting})
\item Lecture : S. Mac Lane et I. Moerdijk, \textit{Sheaves in Geometry and Logic: A First Introduction to Topos Theory}
\end{refs}

\bigskip

\begin{problem}[{Le nerf d'une catégorie --- Master, short}]
\textbf{CAT-32.} \textbf{Problème.} Le \textit{nerf} d'une petite catégorie $ \mathcal{C} $
est l'ensemble simplicial $ N(\mathcal{C}) $ dont les $ n $-simplexes sont les foncteurs
$ [n] \to \mathcal{C} $, où $ [n] $ est l'ensemble ordonné $ 0 \le 1 \le \cdots \le n $ --- de
sorte qu'un $ n $-simplexe est une chaîne de $ n $ morphismes composables. Démontrez que
$ N $ est un foncteur plein et fidèle de $ \mathbf{Cat} $ vers les ensembles simpliciaux, et
qu'un ensemble simplicial est isomorphe à un nerf si et seulement si toute corne interne
$ \Lambda^{n}_{k} \to X $ avec $ 0 < k < n $ admet un remplissage \textit{unique}.
\end{problem}

\noindent Le nerf transforme une catégorie en objet combinatoire et un foncteur en
application simpliciale, de sorte que l'espace classifiant d'une catégorie peut être formé et
toute la machinerie de la théorie de l'homotopie lui être appliquée --- la construction derrière la
K-théorie algébrique supérieure de Quillen et derrière l'identification de la cohomologie des
groupes avec la cohomologie d'un espace. Sa caractérisation ci-dessus est aussi la porte d'entrée
de la théorie des catégories supérieures : supprimez le mot \og{} unique \fg{} de la condition de
remplissage des cornes, en n'exigeant plus que l'existence, et l'on obtient les complexes de Kan
faibles de Boardman et Vogt, rebaptisés quasi-catégories par Joyal, qui sont les $\infty$-catégories de
CAT-13 et CAT-25.

\begin{refs}
\item Contexte : Nerf (théorie des catégories) --- Wikipédia (\url{https://fr.wikipedia.org/wiki/Nerf_%28th%C3%A9orie_des_cat%C3%A9gories%29})
\item Contexte : Ensemble simplicial --- Wikipédia (\url{https://fr.wikipedia.org/wiki/Ensemble_simplicial})
\item Lecture : Emily Riehl, \textit{Categorical Homotopy Theory} (Cambridge, 2014)
\end{refs}

\bigskip

\begin{problem}[{Les transformations naturelles comme une fin --- Master, short}]
\textbf{CAT-33.} \textbf{Problème.} Pour des foncteurs $ F, G \colon \mathcal{C} \to \mathcal{D} $
avec $ \mathcal{C} $ petite et $ \mathcal{D} $ localement petite, démontrez que l'ensemble des
transformations naturelles de $ F $ vers $ G $ est la fin
\[ \mathrm{Nat}(F, G) \;\cong\; \int_{c \in \mathcal{C}} \mathrm{Hom}_{\mathcal{D}}\bigl(F c,\, G c\bigr), \]
c'est-à-dire l'égaliseur des deux applications évidentes
$ \prod_{c} \mathrm{Hom}(Fc, Gc) \rightrightarrows \prod_{u \colon c \to c'} \mathrm{Hom}(Fc, Gc') $.
Déduisez-en la forme \og{} fin \fg{} du lemme de Yoneda (CAT-08) :
$ \int_{c} \mathrm{Hom}\bigl( \mathrm{Hom}(a, c),\, F c \bigr) \cong F(a) $.
\end{problem}

\noindent La construction remonte à l'article de Yoneda de 1960 sur Ext, et Mac Lane lui
a donné la notation intégrale --- justifiée par un théorème de Fubini qui permet d'échanger des fins
itérées exactement comme on échange des intégrales itérées. La formule démontrée ici est la raison
pour laquelle ce calcul mérite d'être appris : les énoncés de naturalité deviennent des égalités
entre intégrales. La cofin duale $ \int^{c} $ fournit l'autre moitié, en réécrivant le théorème
de densité de CAT-20 sous la forme $ X \cong \int^{c} X(c) \cdot \mathrm{Hom}(-, c) $ et en
donnant aux extensions de Kan de CAT-29 des formules closes.

\begin{refs}
\item Contexte : Fin (théorie des catégories) --- Wikipédia (\url{https://fr.wikipedia.org/wiki/Fin_%28th%C3%A9orie_des_cat%C3%A9gories%29})
\item Lecture : Saunders Mac Lane, \textit{Categories for the Working Mathematician}, ch. IX
\item Lecture : Fosco Loregian, \textit{(Co)end Calculus} (Cambridge University Press, 2021)
\end{refs}

\bigskip


\section*{Recherche}

\begin{problem}[{Théorie synthétique et vérifiée par machine des $\infty$-catégories --- Recherche}]
\textbf{CAT-13.} \textbf{Problème.} La théorie classique (\og{} analytique \fg{}) des $\infty$-catégories, telle que développée
par Joyal et Lurie, construit les ($\infty$,1)-catégories à partir des quasi-catégories --- des ensembles
simpliciaux soumis à des conditions de remplissage de cornes --- au prix d'une combinatoire
redoutable. Le programme \textit{synthétique} (Riehl--Shulman, 2017) travaille au contraire dans une
extension dirigée de la théorie homotopique des types dont les objets de base se comportent comme
des $\infty$-catégories, et en 2024 Kudasov, Riehl et Weinberger ont fait vérifier par machine le lemme
de Yoneda $\infty$-catégorique dans l'assistant de preuve Rzk.
\begin{enumerate}
\item[(a)] Comme point d'entrée, formalisez le lemme de Yoneda 1-catégorique (CAT-08) dans un assistant
de preuve de votre choix et comparez votre développement, ligne à ligne, avec le développement
$\infty$-catégorique cité.
\item[(b)] La direction de recherche : déterminez quelle part de la théorie des $\infty$-catégories --- limites et
colimites, l'équivalence straightening/unstraightening, les catégories présentables --- peut être
développée synthétiquement et formalisée, et si les démonstrations synthétiques peuvent être
transportées de façon correcte vers tous les modèles.
\end{enumerate}
\textit{Statut : ouvert et actif ; le lemme de Yoneda et des parties de la théorie élémentaire sont
formalisés, tandis que l'essentiel de la théorie à l'échelle de Lurie ne l'est pas.}
\end{problem}

\noindent La formalisation est passée cette décennie du statut de curiosité à celui de
méthode de travail --- de larges pans des mathématiques modernes vivent désormais dans le Mathlib de
Lean --- et la théorie des catégories en est naturellement une adoptante précoce, puisque ses
démonstrations sont des chasses au diagramme qu'une machine sait vérifier et que ses fondements
mettent la théorie des ensembles à rude épreuve. L'approche synthétique porte un rêve catégorique
classique : que le bon langage rende faciles les théorèmes difficiles, la difficulté étant
remboursée lorsque les résultats s'appliquent d'un coup à tous les modèles.

\begin{refs}
\item Contexte : Théorie des catégories supérieures --- Wikipédia (\url{https://fr.wikipedia.org/wiki/Th%C3%A9orie_des_cat%C3%A9gories_sup%C3%A9rieures})
\item Contexte : Lean (assistant de preuve) --- Wikipédia (\url{https://fr.wikipedia.org/wiki/Lean_%28assistant_de_preuve%29})
\item Article : N. Kudasov, E. Riehl et J. Weinberger, \og{} Formalizing the $\infty$-categorical Yoneda lemma \fg{} (CPP 2024), arXiv:2309.08340 (\url{https://arxiv.org/abs/2309.08340})
\item Article : E. Riehl et M. Shulman, \og{} A type theory for synthetic $\infty$-categories \fg{} (2017), arXiv:1705.07442 (\url{https://arxiv.org/abs/1705.07442})
\end{refs}

\bigskip

\begin{problem}[{Le problème des types semi-simpliciaux --- Recherche}]
\textbf{CAT-14.} \textbf{Problème.} En théorie homotopique des types (HoTT) --- le fondement dans lequel les types
se comportent comme des types d'homotopie et l'égalité comme des chemins --- Voevodsky a demandé,
pendant l'année Univalent Foundations 2012-2013 à l'IAS, si le type des \textit{types
semi-simpliciaux} peut être défini de manière interne : une famille
$ X_0, X_1, X_2, \ldots $ où $ X_0 $ est un type, $ X_1 $ une famille de types au-dessus des
couples issus de $ X_0 $, $ X_2 $ une famille au-dessus des triangles, et ainsi de suite,
toutes les compatibilités de faces valant de manière cohérente.
\begin{enumerate}
\item[(a)] Construisez les trois premiers niveaux à la main à l'intérieur de HoTT.
\item[(b)] Identifiez exactement l'obstruction de cohérence qui bloque une définition uniforme de la tour
infinie.
\end{enumerate}
\textit{Statut : la définissabilité interne des types semi-simpliciaux dans le HoTT du livre est
ouverte --- c'est un problème ouvert central du domaine. Kolomatskaia et Shulman (2023-2024) ont
donné une construction en \og{} théorie des types affichée \fg{}, une extension de HoTT à traits modaux ;
savoir si le HoTT ordinaire suffit reste inconnu.}
\end{problem}

\noindent Le problème a l'air d'une trivialité comptable et n'en est pas une du tout :
écrire \og{} toutes les cohérences supérieures \fg{} est précisément ce qu'un langage dont l'égalité est
homotopique ne sait pas manifestement faire à son propre sujet. Il conditionne l'objectif plus
large de faire de la théorie des catégories supérieures (CAT-13) nativement à l'intérieur de HoTT.
Voevodsky l'a jugé assez sérieux pour concevoir un système étendu (HTS) autour de lui, au prix de
ses propriétés calculatoires ; l'avancée récente de la théorie des types affichée retrouve une
définition tout en conservant un meilleur comportement, au prix d'un fondement non standard.

\begin{refs}
\item Contexte : Théorie homotopique des types --- Wikipédia (\url{https://fr.wikipedia.org/wiki/Th%C3%A9orie_homotopique_des_types})
\item Contexte : Problèmes ouverts en théorie homotopique des types --- nLab (en anglais) (\url{https://ncatlab.org/nlab/show/open+problems+in+homotopy+type+theory})
\item Article : A. Kolomatskaia et M. Shulman, \og{} Displayed type theory and semi-simplicial types \fg{} (2023), arXiv:2311.18781 (\url{https://arxiv.org/abs/2311.18781})
\item Lecture : The Univalent Foundations Program, \textit{Homotopy Type Theory: Univalent Foundations of Mathematics} (2013, librement accessible sur homotopytypetheory.org/book (\url{https://homotopytypetheory.org/book/}))
\end{refs}

\bigskip

\begin{problem}[{Univers univalents dans tout ($\infty$,1)-topos --- Recherche}]
\textbf{CAT-15.} \textbf{Problème.} L'axiome d'univalence de Voevodsky affirme que pour des types
$ A, B $ d'un univers, l'application canonique du type d'identité $ A = B $ vers le type
des équivalences $ A \simeq B $ est elle-même une équivalence --- \og{} les choses isomorphes sont
égales \fg{}. La conjecture sémantique (Awodey, Joyal et d'autres, au début des années 2010) était que
la théorie homotopique des types à univers univalents admet des modèles dans tout
($\infty$,1)-topos de Grothendieck, faisant de HoTT un langage interne pour la théorie des topos
supérieurs.
\begin{enumerate}
\item[(a)] Traitez le premier cas non trivial : dans le modèle groupoïdal de Hofmann--Streicher de la
théorie des types, vérifiez que l'univers des ensembles (les groupoïdes discrets) satisfait
l'univalence.
\item[(b)] Étudiez l'énoncé général et sa démonstration.
\textit{Statut : résolu pour les ($\infty$,1)-topos de Grothendieck --- démontré par Shulman en 2019 au moyen de
structures de modèles injectives et d'univers univalents stricts. L'énoncé analogue pour les
($\infty$,1)-topos élémentaires reste ouvert, ainsi que la recherche d'une définition pleinement
satisfaisante d'\og{} ($\infty$,1)-topos élémentaire \fg{} elle-même.}
\end{enumerate}

\end{problem}

\noindent Ce théorème est le mur porteur du programme univalent : il garantit qu'un
théorème démontré dans HoTT est simultanément un théorème sur les espaces, sur les faisceaux
d'espaces, sur les espaces équivariants et sur tout autre ($\infty$,1)-topos de Grothendieck. Le modèle
groupoïdal de 1998 de Hofmann et Streicher, construit bien avant que l'univalence ne soit nommée,
contient l'idée entière en miniature --- c'est pourquoi la partie (a) est une préparation honnête à
la partie (b), une démonstration qui traverse de la sérieuse théorie des catégories de modèles.

\begin{refs}
\item Contexte : Fondements univalents --- Wikipédia (\url{https://fr.wikipedia.org/wiki/Fondements_univalents})
\item Article : M. Shulman, \og{} All ($\infty$,1)-toposes have strict univalent universes \fg{} (2019), arXiv:1904.07004 (\url{https://arxiv.org/abs/1904.07004})
\item Lecture : The Univalent Foundations Program, \textit{Homotopy Type Theory: Univalent Foundations of Mathematics} (2013)
\end{refs}

\bigskip

\begin{problem}[{L'hypothèse d'homotopie de Grothendieck --- Recherche, short}]
\textbf{CAT-23.} \textbf{Problème.} Dans \textit{Pursuing Stacks} (1983), Grothendieck a donné une
définition purement algébrique de l'$\infty$-groupoïde faible et a prédit que la construction de
l'$\infty$-groupoïde fondamental induit une équivalence entre la théorie de l'homotopie des espaces
topologiques et la théorie de l'homotopie des $\infty$-groupoïdes de Grothendieck. Démontrez ou réfutez
cette hypothèse d'homotopie pour la définition propre à Grothendieck.
\end{problem}

\noindent \textbf{Statut : ouvert} pour les $\infty$-groupoïdes algébriques de Grothendieck. Pour
des modèles combinatoires comme les complexes de Kan, la correspondance est vraie essentiellement
par construction, de sorte que le contenu du problème est de savoir si un $\infty$-groupoïde défini
algébriquement --- opérations de composition et cohérences engendrées par une théorie globulaire ---
capture exactement le type d'homotopie d'un espace. Kapranov et Voevodsky ont publié en 1991 une
démonstration d'une version de l'énoncé ; l'argument s'est révélé plus tard erroné, et la forme
affinée de ce qu'ils affirmaient est aujourd'hui la conjecture de Simpson de CAT-24. C'est
l'article de foi fondateur de la théorie des catégories supérieures, et c'est encore une
conjecture.

\begin{refs}
\item Contexte : Hypothèse d'homotopie --- Wikipédia (en anglais) (\url{https://en.wikipedia.org/wiki/Homotopy_hypothesis})
\item Contexte : $\infty$-groupoïde --- Wikipédia (en anglais) (\url{https://en.wikipedia.org/wiki/Infinity-groupoid})
\item Contexte : Hypothèse d'homotopie --- nLab (en anglais) (\url{https://ncatlab.org/nlab/show/homotopy+hypothesis})
\end{refs}

\bigskip

\begin{problem}[{À quel point une \textit{n}-catégorie faible peut-elle être stricte ? --- Recherche, short}]
\textbf{CAT-24.} \textbf{Problème.} Gordon, Power et Street ont démontré en 1995 que toute
tricatégorie est tri-équivalente à une Gray-catégorie --- associativité stricte, échange faible ---
mais non, en général, à une 3-catégorie stricte. Tranchez la question analogue de
semi-strictification en dimension quatre et au-delà : la conjecture de Simpson, selon laquelle
toute $ n $-catégorie faible est équivalente à une autre dans laquelle tout est
strict sauf les unités.
\end{problem}

\noindent \textbf{Statut : théorème en dimension trois, ouvert en général.} Simpson a montré
en 1998 que les 3-groupoïdes stricts ne peuvent pas modéliser tous les 3-types d'homotopie --- le
produit de Whitehead sur $ S^2 $ est l'obstruction --- de sorte qu'une part de faiblesse est
réellement nécessaire, et sa conjecture en isole le minimum : gardez les unités faibles,
strictifiez tout le reste. La question est le visage pratique de l'hypothèse d'homotopie
(CAT-23) : toute définition praticable de $ n $-catégorie faible paie la cohérence quelque part,
et la semi-strictification demande à quel point ce paiement peut être réduit.

\begin{refs}
\item Contexte : 3-catégorie (tricatégories) --- Wikipédia (en anglais) (\url{https://en.wikipedia.org/wiki/3-category})
\item Contexte : Théorie des catégories supérieures --- Wikipédia (\url{https://fr.wikipedia.org/wiki/Th%C3%A9orie_des_cat%C3%A9gories_sup%C3%A9rieures})
\item Lecture : R. Gordon, A. J. Power et R. Street, \og{} Coherence for tricategories \fg{} (1995), Memoirs of the American Mathematical Society 117
\item Article : C. Simpson, \og{} Homotopy types of strict 3-groupoids \fg{} (1998), arXiv:math/9810059 (\url{https://arxiv.org/abs/math/9810059})
\end{refs}

\bigskip

\begin{problem}[{L'hypothèse du cobordisme --- Recherche}]
\textbf{CAT-25.} \textbf{Problème.} Une théorie topologique des champs de dimension $ n $ est un foncteur
monoïdal symétrique défini sur une catégorie de bordismes ; la version \textit{étendue} permet de
découper les variétés jusqu'aux points, produisant une $ (\infty, n) $-catégorie monoïdale
symétrique $ \mathrm{Bord}_n^{\mathrm{fr}} $ de bordismes framés. Baez et Dolan ont conjecturé en
1995 que cette catégorie est la $ (\infty, n) $-catégorie monoïdale symétrique libre avec duaux
sur un unique générateur --- de sorte qu'une théorie des champs framée étendue est déterminée par
l'unique objet qu'elle attribue à un point, lequel peut être n'importe quel objet
\textit{pleinement dualisable} de la cible.
\begin{enumerate}
\item[(a)] Traitez les cas de basse dimension à la main : démontrez qu'un foncteur monoïdal symétrique
$ \mathrm{Bord}_1 \to \mathrm{Vect}_k $ force l'espace vectoriel attaché à un point à être de
dimension finie, les bordismes en zigzag réalisant la dualité, et que la théorie en est
déterminée.
\item[(b)] Démontrez la classification des théories topologiques des champs (non étendues) de dimension
deux : elles correspondent aux algèbres de Frobenius commutatives, le pantalon fournissant la
multiplication et le disque la trace.
\item[(c)] Énoncez précisément l'hypothèse du cobordisme comme une équivalence entre l'espace des théories
des champs et l'espace des objets pleinement dualisables, ainsi que sa version pour les variétés à
$ G $-structure tangentielle, et parcourez la stratégie de démonstration : théorie de Morse et
décompositions en anses du côté géométrique, données de dualisabilité du côté algébrique.
\end{enumerate}
\textit{Statut : largement admise, mais sa vérification écrite est encore en train de se stabiliser.
Lurie a publié en 2009 une esquisse détaillée que le domaine utilise quotidiennement ; des
démonstrations complètes suivant d'autres voies ont été proposées depuis (Ayala--Francis, 2017 ;
Grady--Pavlov, 2021), et tous les détails de l'approche propre à Lurie ne sont pas parus. Traitez-la
comme un théorème dont la comptabilité reste une tâche de recherche vivante.}
\end{problem}

\noindent C'est ici que l'abstraction de cette page touche la physique. Atiyah et Segal
ont axiomatisé les théories des champs comme des foncteurs à la fin des années 1980 ; Baez et Dolan
ont ensuite proposé que les variétés elles-mêmes forment une catégorie supérieure libre, si bien
qu'une théorie des champs étendue tout entière se comprime en un unique objet dualisable --- une
version catégorifiée du fait qu'une représentation d'un groupe libre n'est qu'un choix de matrices.
L'énoncé est aussi la meilleure publicité existante pour la théorie des catégories supérieures : il
dit qu'une question subtile sur les variétés lisses, en toutes dimensions, se règle par une
propriété universelle.

\begin{refs}
\item Contexte : Hypothèse du cobordisme --- Wikipédia (en anglais) (\url{https://en.wikipedia.org/wiki/Cobordism_hypothesis})
\item Contexte : Théorie quantique topologique des champs --- Wikipédia (en anglais) (\url{https://en.wikipedia.org/wiki/Topological_quantum_field_theory})
\item Article : J. Baez et J. Dolan, \og{} Higher-dimensional algebra and topological quantum field theory \fg{} (1995), J. Math. Phys. 36
\item Article : J. Lurie, \og{} On the classification of topological field theories \fg{} (2009), arXiv:0905.0465 (\url{https://arxiv.org/abs/0905.0465})
\item Article : D. Ayala et J. Francis, \og{} The cobordism hypothesis \fg{} (2017), arXiv:1705.02240 (\url{https://arxiv.org/abs/1705.02240})
\item Lecture : Joachim Kock, \textit{Frobenius Algebras and 2D Topological Quantum Field Theories} (Cambridge University Press, LMS Student Texts 59)
\end{refs}

\bigskip

\begin{problem}[{Catégories triangulées sans modèle --- Recherche}]
\textbf{CAT-34.} \textbf{Problème.} Une \textit{catégorie triangulée} est une catégorie additive munie d'un
automorphisme de décalage $ X \mapsto X[1] $ et d'une classe de triangles distingués
$ X \to Y \to Z \to X[1] $ satisfaisant les axiomes de Verdier ; la catégorie dérivée d'un
anneau ou d'un espace en est l'exemple motivant. Les axiomes sont notoirement déficients sur un
point : le troisième sommet $ Z $ d'un triangle bâti sur un morphisme donné n'est déterminé qu'à
isomorphisme non canonique près.
\begin{enumerate}
\item[(a)] Rendez le défaut précis. Démontrez à partir des axiomes que le cône d'un morphisme est unique à
isomorphisme près, puis exhibez un carré commutatif dont le morphisme de cônes induit n'est pas
déterminé de manière unique --- de sorte que \og{} cône \fg{} ne peut pas être rendu fonctoriel sur la
catégorie des flèches.
\item[(b)] Comprenez la réparation. Un \textit{enrichissement} d'une catégorie triangulée est une
dg-catégorie ou une $\infty$-catégorie stable dont elle est la catégorie homotopique ; dans un
enrichissement, les cônes sont d'authentiques colimites et sont donc fonctoriels. Vérifiez que la
catégorie homotopique d'une $\infty$-catégorie stable est triangulée, et identifiez lequel des axiomes de
Verdier --- l'axiome de l'octaèdre en particulier --- devient un théorème plutôt qu'une hypothèse.
\item[(c)] Lisez les deux théorèmes négatifs et localisez l'obstruction que chacun exploite. Muro,
Schwede et Strickland (2007) ont construit une catégorie triangulée qui n'est la catégorie
homotopique d'aucune catégorie de modèles stable ; Rizzardo et Van den Bergh (2018) en ont
construit une qui est de plus $ k $-linéaire sur un corps, le cas que l'exemple antérieur
n'atteignait pas, au moyen d'une théorie de l'obstruction bâtie sur la cohomologie de Hochschild.
\item[(d)] La question de recherche : caractérisez les catégories triangulées qui admettent un
enrichissement, et celles dont l'enrichissement est unique. Lunts et Orlov (J. Amer. Math. Soc. 23,
2010) ont démontré l'unicité pour les catégories dérivées d'une large classe de schémas, et l'on
connaît des exemples à enrichissement non unique ; aucun critère général n'existe.
\end{enumerate}
\textit{Statut : (a)--(c) relèvent de la théorie établie ; (d) est ouvert. L'existence et l'unicité
des enrichissements sont établies dans de nombreux cas particuliers et fausses dans d'autres, sans
aucune caractérisation générale connue.}
\end{problem}

\noindent Jean-Louis Verdier a introduit les catégories triangulées dans sa thèse de 1963
sous la direction de Grothendieck, pour donner à la catégorie dérivée d'une catégorie abélienne une
vie axiomatique propre. Les axiomes consignent ce qui survit d'une théorie de l'homotopie une fois
passé aux classes d'homotopie d'applications ; la partie (a) dit ce qui est perdu. Pendant quarante
ans on ignora si cette perte était jamais fatale --- si toute catégorie triangulée est l'ombre de
quelque chose de mieux élevé. Muro, Schwede et Strickland ont répondu non en 2007, et Rizzardo et
Van den Bergh ont levé le doute restant pour les catégories $ k $-linéaires en 2018. La réponse
pratique a été de cesser de travailler dans l'ombre : dg-catégories et $\infty$-catégories stables, où un
cône est une colimite et la fonctorialité est gratuite, ce qui est la principale raison pour
laquelle les $\infty$-catégories sont devenues la langue de travail de l'algèbre homologique.

\begin{refs}
\item Contexte : Catégorie triangulée --- Wikipédia (\url{https://fr.wikipedia.org/wiki/Cat%C3%A9gorie_triangul%C3%A9e})
\item Contexte : Catégorie dérivée --- Wikipédia (\url{https://fr.wikipedia.org/wiki/Cat%C3%A9gorie_d%C3%A9riv%C3%A9e})
\item Article : F. Muro, S. Schwede et N. Strickland, \og{} Triangulated categories without models \fg{}, Inventiones Mathematicae 170 (2007), arXiv:0704.1378 (\url{https://arxiv.org/abs/0704.1378})
\item Article : A. Rizzardo et M. Van den Bergh, \og{} A k-linear triangulated category without a model \fg{} (2018), arXiv:1801.06344 (\url{https://arxiv.org/abs/1801.06344})
\end{refs}

\bigskip

\begin{problem}[{Localisation, orthogonalité et grands cardinaux --- Recherche, short}]
\textbf{CAT-35.} \textbf{Problème.} Le principe de Vopěnka --- l'axiome de grand cardinal
affirmant que dans toute classe propre de graphes il existe deux membres distincts reliés par un
homomorphisme de l'un vers l'autre --- entraîne à la fois que toute sous-catégorie pleine d'une
catégorie localement présentable stable par limites est réflexive, et que des foncteurs de
localisation homotopique existent pour toute théorie cohomologique. Déterminez lesquelles de ces
conséquences sont démontrables dans ZFC seul, et cernez la force exacte en grands cardinaux de
chacune.
\end{problem}

\noindent \textbf{Statut : partiellement réglé, partiellement ouvert.} Adámek et Rosický ont
fait de ce cercle d'idées un chapitre de \textit{Locally Presentable and Accessible Categories}
(1994) : réflexivité des sous-catégories stables par limites, petitesse des classes
d'orthogonalité et principe de Vopěnka se révèlent être presque le même énoncé, de sorte qu'une
question toute simple sur l'existence d'une construction devient une question sur la taille de
l'univers ensembliste. On sait que l'énoncé de réflexivité porte une force de grand cardinal et
n'est donc pas démontrable dans ZFC ; la force exacte de la forme faible du principe de Vopěnka
fait l'objet de travaux continus. Casacuberta, Scevenels et Smith ont porté le lien en théorie de
l'homotopie en 2005, en montrant que sous le principe de Vopěnka tout foncteur homotopiquement
idempotent sur les espaces est une localisation par rapport à une unique application --- alors que
l'énoncé correspondant pour les spectres avait été démontré sans détour par Bousfield en 1979,
sans aucun grand cardinal. Bagaria, Casacuberta, Mathias et Rosický ont montré plus tard que pour
les classes \textit{définissables} les cardinaux requis chutent radicalement. La théorie des
catégories rencontre rarement la théorie des ensembles aussi directement : ici, la réponse à
\og{} cette localisation existe-t-elle ? \fg{} dépend des axiomes sur lesquels on se tient.

\begin{refs}
\item Contexte : Principe de Vopěnka --- Wikipédia (en anglais) (\url{https://en.wikipedia.org/wiki/Vop%C4%9Bnka%27s_principle})
\item Contexte : Catégorie accessible --- Wikipédia (en anglais) (\url{https://en.wikipedia.org/wiki/Accessible_category})
\item Lecture : J. Adámek et J. Rosický, \textit{Locally Presentable and Accessible Categories} (Cambridge, LMS Lecture Note Series 189, 1994), ch. 6
\item Article : C. Casacuberta, D. Scevenels et J. H. Smith, \og{} Implications of large-cardinal principles in homotopical localization \fg{}, Advances in Mathematics 197 (2005), 120--139
\end{refs}

\bigskip


\section*{Pour aller plus loin}


\vfill
\noindent\emph{Hors de la Caverne} --- des probl\`emes, pas de r\'eponses.
\end{document}
